\documentclass[11pt,a4paper]{article}

% ---- Encoding & fonts ----
\usepackage[T1]{fontenc}
\usepackage[utf8]{inputenc}
\usepackage{lmodern}

% ---- Layout ----
\usepackage[margin=1in]{geometry}
\usepackage{setspace}
\onehalfspacing

% ---- Math & tables ----
\usepackage{amsmath,amssymb,amsthm}
\usepackage{booktabs}
\usepackage{multirow}
\usepackage{tabularx}
\usepackage{array}

% ---- Figures & graphics ----
\usepackage{graphicx}
\usepackage{xcolor}

% ---- Lists & typography ----
\usepackage{enumitem}
\usepackage{microtype}
\usepackage{hyperref}
\hypersetup{
  colorlinks=true,
  linkcolor=blue!50!black,
  citecolor=blue!50!black,
  urlcolor=blue!60!black
}

% ---- Bibliography ----
\usepackage[numbers,sort&compress]{natbib}
\bibliographystyle{plainnat}

% ---- Theorem-like environments ----
\theoremstyle{definition}
\newtheorem{definition}{Definition}
\newtheorem{remark}{Remark}

\title{%
  Randomness Quality in Machine Learning and Stochastic Optimisation:\\
  A Critical Literature Review and Research Agenda
}
\author{%
  % Authors to be filled
  Anonymous / Project Randomness
}
\date{\today}

\begin{document}
\maketitle

\begin{abstract}
Random numbers underpin nearly every stochastic component of modern machine learning (ML) and optimisation: weight initialisation, mini-batch sampling, dropout and other stochastic regularisers, data augmentation, evolutionary operators, and the noise that shapes the trajectory of stochastic gradient methods.
Despite this centrality, the \emph{statistical quality} of the underlying random bitstreams---as opposed to the mere presence of a seed---remains inconsistently measured, inconsistently controlled, and inconsistently reported.
This paper provides a critical review of the literature at the intersection of random number generation (RNG) quality, deep learning training pipelines, and stochastic / metaheuristic optimisation.
We synthesise evidence from (i)~classical and quantum RNG characterisation, (ii)~empirical studies that substitute pseudo-random, true-random, and quantum-random sources into neural training, (iii)~theory and empirics on how gradient noise structures the solutions found by SGD and adaptive optimisers, and (iv)~metaheuristic studies of RNG quality in genetic algorithms, particle swarm optimisation, and differential evolution.
Our critique highlights contradictory findings (notably on quantum versus classical initialisation), pervasive underpowering and confounded experimental designs, and a persistent gap between \emph{cryptographic / statistical} notions of randomness quality and the \emph{geometry-aware} noise that matters for optimiser dynamics.
We conclude with a structured research agenda for comparing large binary random archives in AI/ML tasks under isolated stochastic sites, multi-seed statistics, and explicit quality tiers.
\end{abstract}

\tableofcontents
\newpage

\section{Introduction}
\label{sec:intro}

\subsection{Motivation}

Machine learning systems are stochastic artefacts.
Even when the architecture, loss, and hyperparameters are fixed, training outcomes vary with the realisation of random draws used for initialisation, data ordering, regularisation masks, and augmentation~\cite{summers2021nondeterminism,zhuang2022randomness,pham2022quantifying}.
Optimisation methods themselves---stochastic gradient descent (SGD) and its adaptive variants, as well as population-based metaheuristics---are defined by how they inject and structure randomness into the search process~\cite{keskar2017largebatch,smith2018bayesian,zhu2019anisotropic}.

A separate, older literature asks a different question: when is a sequence of bits ``random enough''?
Statistical test batteries such as NIST~SP~800-22, Dieharder, and TestU01~\cite{bassham2010nist,lecuyer2007testu01} and, more recently, neural distinguishers~\cite{valle2024assessing} attempt to answer that question for cryptography and simulation.
The present review sits at the junction of these two traditions.
If an organisation holds \emph{very large binary random archives}---from high-quality PRNGs, hardware TRNGs, or QRNGs---can differences in measured bitstream quality be shown to matter for AI/ML task quality, training stability, or optimiser behaviour?

\subsection{Scope}

We review and critically assess work on:
\begin{enumerate}[leftmargin=*]
  \item definitions and measurement of RNG quality (classical batteries; entropy estimates; ML-based tests);
  \item substitution of PRNG / TRNG / QRNG sources into neural network training (initialisation, dropout, shuffling);
  \item the role of randomness and noise structure in first-order optimisers (SGD, Adam and relatives) and in metaheuristics (GA, PSO, DE);
  \item reproducibility, tooling nondeterminism, and reporting standards for seed-sensitive studies.
\end{enumerate}
We deliberately treat \emph{seed sensitivity} (different seeds of the same generator) and \emph{generator quality} (different statistical properties of the source) as distinct constructs; conflating them is a recurring methodological failure in the literature (\S\ref{sec:critique}).

\subsection{Contributions of this review}

\begin{enumerate}[leftmargin=*]
  \item A structured map of how randomness enters ML and optimisation pipelines, and which notions of ``quality'' are claimed to matter.
  \item A critical synthesis of conflicting empirical claims (especially QRNG vs.\ PRNG), with attention to statistical power, effect size, and confounds.
  \item An explicit gap analysis connecting cryptographic RNG evaluation, framework PRNG implementations, and optimiser noise theory.
  \item A research agenda for controlled experiments with large binary random bitstreams (\S\ref{sec:agenda}), aligned with the accompanying project plan.
\end{enumerate}

\subsection{Organisation}

Section~\ref{sec:background} fixes terminology and quality metrics.
Section~\ref{sec:rngml} reviews RNG substitution studies in ML.
Section~\ref{sec:opt} reviews randomness in stochastic and metaheuristic optimisation.
Section~\ref{sec:critique} offers a critical synthesis.
Section~\ref{sec:gaps} states research gaps.
Section~\ref{sec:agenda} outlines a concrete experimental programme.

\section{Background: Randomness, Quality, and Where It Enters Learning Systems}
\label{sec:background}

\subsection{Sources of randomness}

\begin{definition}[PRNG]
A \emph{pseudorandom number generator} is a deterministic algorithm that expands a short seed into a long sequence intended to be computationally indistinguishable from i.i.d.\ uniform draws for a target class of statistical tests or adversaries.
\end{definition}

\begin{definition}[TRNG / physical RNG]
A \emph{true} (physical) random number generator harvests entropy from a classical physical process (thermal noise, jitter, etc.), typically followed by conditioning / extraction.
\end{definition}

\begin{definition}[QRNG]
A \emph{quantum} random number generator harvests unpredictability justified by quantum measurement (e.g., photon arrival, vacuum fluctuations, phase noise)~\cite{herrero2017quantum,mannalath2022qrng}.
\end{definition}

In ML practice, the dominant sources are software PRNGs shipped with NumPy, PyTorch, and TensorFlow (historically Mersenne Twister; increasingly Philox / PCG-class counter-based generators)~\cite{antunes2025statistical,antunes2024reproducibility}.
Hardware TRNG/QRNG streams appear mainly in research substitutions~\cite{bird2019softcomputing,koivu2022dropout,heese2024biased}.

\subsection{What ``quality'' means}

The literature uses several non-equivalent notions:

\begin{enumerate}[leftmargin=*]
  \item \textbf{Uniformity and independence tests.}
  Batteries such as NIST~SP~800-22, Dieharder, and TestU01 SmallCrush/Crush/BigCrush~\cite{bassham2010nist,lecuyer2007testu01} report $p$-values against specific alternatives (bias, runs, serial correlation, linear complexity, etc.).
  Passing is necessary but not sufficient for cryptographic use~\cite{soto1999statistical}.

  \item \textbf{Entropy / min-entropy.}
  Especially for physical sources after digitisation; post-processing aims to extract near-uniform bits~\cite{herrero2017quantum}.

  \item \textbf{Predictability / next-bit hardness.}
  Cryptographic security notions; relevant when adversaries may recover seeds~\cite{dahiya2024mlrandomness}.

  \item \textbf{ML distinguishability.}
  Neural networks trained to classify windows of bits against a ``golden'' reference~\cite{valle2024assessing}; longer windows and CNNs tend to be more discriminative for weak PRNGs.

  \item \textbf{Task-level utility.}
  Downstream accuracy, calibration, or convergence speed when the stream is injected into an ML or optimisation algorithm~\cite{bird2019softcomputing,koivu2022dropout,bastos2009pso}.
\end{enumerate}

\begin{remark}
These notions can disagree.
A stream that fails Crush may still train an MNIST MLP indistinguishably from \texttt{os.urandom}; conversely, a stream that passes NIST may still be a poor match to the \emph{anisotropic} noise SGD ``wants'' for flat-minima selection (\S\ref{sec:opt}).
Conflating cryptographic quality with optimisation utility is a central conceptual gap (\S\ref{sec:gaps}).
\end{remark}

\subsection{Stochastic sites in modern ML pipelines}

Following tooling analyses~\cite{zhuang2022randomness,summers2021nondeterminism}, we partition randomness into sites that can be ablated independently:

\begin{table}[t]
\centering
\caption{Primary stochastic sites in neural training pipelines.}
\label{tab:sites}
\begin{tabular}{@{}llp{6.2cm}@{}}
\toprule
\textbf{ID} & \textbf{Site} & \textbf{Role of randomness} \\
\midrule
S1 & Weight / bias initialisation & Sets the starting point on a non-convex loss; interacts with activation scale (Xavier/He) \\
S2 & Stochastic regularisation & Dropout, DropConnect, stochastic depth masks \\
S3 & Data order / mini-batches & Induces gradient noise in SGD-family methods \\
S4 & Data augmentation & Random crops, flips, colour jitter, MixUp, etc. \\
S5 & Optimiser internals & Rare explicit noise (e.g., Langevin); Adam's early-step variance is data-dependent, not RNG-seeded in the usual sense~\cite{liu2020radam} \\
S6 & Evaluation / decoding & Temperature sampling, stochastic beams (inference-time) \\
\bottomrule
\end{tabular}
\end{table}

Population-based methods add further sites: population initialisation, crossover/mutation draws, and velocity/position perturbations in PSO~\cite{bastos2009pso,cantu2008rngga}.

\subsection{Seed sensitivity versus generator quality}

Changing the \emph{seed} of a fixed PRNG re-samples the same distributional family.
Changing the \emph{generator} (or feeding bits from a physical archive) can change bias, correlation structure, period, and lattice artefacts.
Most reproducibility papers study the former~\cite{picard2021torch,bethard2022seeds}; only a thinner literature studies the latter~\cite{bird2019softcomputing,koivu2022dropout,heese2024biased,antunes2025influence}.
Both matter, but they answer different scientific questions.
A valid RNG-quality study must hold the experimental design fixed while varying measured quality---not merely re-seeding.

\section{Random Number Quality in Machine Learning: Empirical Landscape}
\label{sec:rngml}

\subsection{Weight initialisation: QRNG versus PRNG}

Bird et al.~\cite{bird2019softcomputing} compared CPU PRNG and QPU-derived QRNG initialisation across dense and convolutional networks.
They report task-dependent gaps: negligible on MNIST ($\approx 0.02\%$), larger on CIFAR-10 ($\approx +0.92\%$ for QRNG) and mental-state EEG ($\approx +2.82\%$), with qualitatively different learning curves.
The authors conclude that the two sources are ``inexplicably'' distinguishable in soft computing practice despite theoretical expectations that high-quality PRNGs should be interchangeable with true randomness for these scales.

Heese et al.~\cite{heese2024biased} directly challenge the narrative of systematic QRNG advantage.
Using hardware-biased quantum random numbers from real devices, unbiased QRNG, classical PRNG, and a classical generator engineered to match hardware bias, they find \emph{no statistically significant} effect of RNG choice on CNN and RNN initialisation.
They explicitly note failure to replicate Bird et al.'s claim under different but ``typical'' neural setups.

\begin{remark}[Critical reading]
The Bird--Heese tension is the clearest illustration that the field lacks standardised protocols.
Neither line of work fully isolates quality metrics of the bitstreams (NIST/TestU01 profiles) as covariates; sample sizes are modest relative to the small claimed effect sizes; and data-split randomness often co-varies with initialisation.
Until studies jointly report (a)~bitstream quality scores, (b)~fixed splits, (c)~$\geq 25$ seeds for sub-$2.5\%$ effects~\cite{akesson2024reporting}, claims of QRNG superiority remain provisional.
\end{remark}

\subsection{Dropout and stochastic regularisation}

Koivu et al.~\cite{koivu2022dropout} replace the RNG driving node dropout with generators of varying statistical quality, including true random sources, across four classification tasks.
Their central finding is nuanced: true/high-quality randomness can be \emph{advantageous or disadvantageous} for overfitting, depending on dataset and hyper-parameters.
This is one of the few studies that treats RNG quality as a first-class experimental factor for a site other than initialisation.

Related work on contextual neural networks also reports sensitivity to the choice of generator during training~\cite{huk2021contextual}.
However, mechanism remains unclear: is the effect mediated by mask correlation structure, effective regularisation strength, or finite-sample artefacts of short random buffers?

\subsection{Framework PRNGs: statistical quality and reproducibility}

Antunes and collaborators~\cite{antunes2025statistical,antunes2024reproducibility} stress-test the PRNGs embedded in NumPy, PyTorch, and TensorFlow against reference C implementations using TestU01 BigCrush on hundreds of streams.
Key findings:
\begin{itemize}[leftmargin=*]
  \item Generators marketed as ``crush-resistant'' (PCG, Philox) can still fail individual tests under some configurations.
  \item Framework integrations do not always match native failure profiles---implementation and seeding (seed$\to$state maps) matter.
  \item Cross-framework numerical reproducibility fails even for ``the same'' algorithm and seed, because seeding functions differ.
\end{itemize}
A follow-on study on the influence of RNG quality in stochastic simulation and ML~\cite{antunes2025influence} argues that, \emph{above a robustness threshold}, PRNG choice is often dominated by performance and convenience, while \emph{low-quality} generators can inject systematic error.
This ``threshold'' view is important: it predicts flat task metrics among strong generators and sharp degradation only for weak/biased ones---a hypothesis our proposed experiments are designed to test.

\subsection{Neural networks as RNG assessors}

Valle et al.~\cite{valle2024assessing} reverse the usual direction: they train CNNs and LSTMs to distinguish candidate RNG outputs from a golden-standard RNG.
Simple LCGs and RC4-derived streams are often distinguishable; stronger constructions and some VCSEL-based QRNG streams are not, even with relatively large networks.
Longer sequence windows and CNN architectures improve discrimination.
This line of work (see also earlier ANN-based CSPRNGs tests~\cite{schindler2018testing}) suggests a complementary Stage-1 metric for ML experiments: if a neural distinguisher cannot tell stream $A$ from stream $B$, expecting large Stage-2 task differences becomes less plausible \emph{a priori}.

Li and Zeng~\cite{li2026transformers} show that Transformers can simulate classical PRNGs (LCG, Mersenne Twister) and generate sequences that pass most NIST tests, while remaining subject to prediction-attack evaluation.
This blurs the boundary between ``ML using RNGs'' and ``ML as RNGs,'' but does not resolve whether NIST-passing streams differentially affect training.

\subsection{Seed-induced variance without changing the generator}

A large body of work documents that \emph{seeds alone} induce non-trivial variance:
\begin{itemize}[leftmargin=*]
  \item Summers and Dinneen~\cite{summers2021nondeterminism} show extreme optimisation instability: perturbing a single weight by machine epsilon can rival other nondeterminism sources.
  \item Zhuang et al.~\cite{zhuang2022randomness} separate algorithmic randomness from implementation/tooling noise (cuDNN, hardware).
  \item Pham et al.~\cite{pham2022quantifying} quantify contributions from initialisation, shuffling, dropout, and splits; data splitting often dominates.
  \item Recent LLM fine-tuning studies show both macro (metric) and micro (per-example) seed effects~\cite{li2025macroseeds}.
\end{itemize}
These results set a floor on detectable RNG-quality effects: any quality-driven delta must exceed seed-driven variance under the same evaluation protocol, or be estimated with adequate replication.

\subsection{Security angle}

Dahiya et al.~\cite{dahiya2024mlrandomness} argue that ML needs stronger randomness standards, showing PRNG-related attack surfaces in contexts such as randomized smoothing.
This is orthogonal to accuracy but relevant if large shared bit archives are reused across experiments without cryptographic hygiene.

\subsection{Interim summary}

Empirically, the ML literature supports three soft conclusions:
\begin{enumerate}[leftmargin=*]
  \item Weak or biased randomness can matter (especially as a negative control).
  \item Among strong PRNGs/QRNGs, task-level differences are small, unstable across papers, and often not significant under careful tests~\cite{heese2024biased,antunes2025influence}.
  \item Site of injection (init vs.\ dropout vs.\ shuffle) is under-explored relative to generator brand names.
\end{enumerate}

\section{Randomness and Optimisation Methods}
\label{sec:opt}

Optimisation is where randomness stops being a mere implementation detail and becomes part of the algorithm's inductive bias.
We separate \emph{first-order stochastic optimisation} (SGD family) from \emph{population-based metaheuristics}, then discuss adaptive methods and open coupling questions.

\subsection{SGD noise, flat minima, and generalisation}

\subsubsection{Batch noise as implicit regularisation}

Keskar et al.~\cite{keskar2017largebatch} popularised the observation that large-batch training often generalises worse than small-batch training, correlating the gap with convergence to sharper minima.
The standard interpretation is that mini-batch sampling injects noise into gradient estimates; smaller batches increase that noise and favour flatter regions.

Smith and Le~\cite{smith2018bayesian} cast this in a Bayesian light, identifying a noise scale
\begin{equation}
  g \approx \varepsilon \frac{N}{B},
\end{equation}
with learning rate $\varepsilon$, dataset size $N$, and batch size $B$.
Holding $\varepsilon$ fixed, there exists an intermediate batch size that maximises test performance---too little noise (large $B$) harms generalisation; too much can harm optimisation.

Zhu et al.~\cite{zhu2019anisotropic} emphasise that SGD noise is \emph{anisotropic}: its covariance structure, not only its scalar magnitude, governs escape from sharp minima.
Wu et al.~\cite{wu2022alignment} prove a stability viewpoint: for over-parameterised models with square loss, linear stability of a minimum under SGD constrains the Frobenius norm of the Hessian in terms of batch size and learning rate, exploiting an alignment property whereby noise concentrates along sharp directions and scales with loss.
Related work on dynamical stability further links large learning rates to flatness-seeking behaviour~\cite{wu2023implicit}.

\begin{remark}[Implication for RNG quality]
Almost all of this theory models gradient noise as arising from \emph{uniform sampling of data} (or with replacement approximations), not from defects in the PRNG that implements the shuffle.
If the shuffle RNG is biased or low-period, the induced noise covariance could deviate from the assumed structure---a theoretically motivated but largely untested pathway from bitstream quality to optimiser geometry.
\end{remark}

\subsubsection{SGD versus adaptive methods}

Adaptive methods (Adam, RMSProp, etc.) reshape effective step sizes per coordinate.
Zhou et al.~\cite{zhou2020adamsgd} analyse why SGD often generalises better than Adam, relating escape behaviour and noise tails to flatter minima.
Liu et al.~\cite{liu2020radam} show that Adam's early adaptive learning-rate estimates have undesirably large variance when few samples have been seen, motivating warmup and Rectified Adam (RAdam).
Here the ``randomness'' is primarily \emph{data-order stochasticity} feeding moment estimates, again mediated by shuffling RNGs and batch construction.

Empirically, comparisons of optimisers that ignore seed variance are unreliable: optimiser rankings can flip across seeds when effect sizes are small~\cite{akesson2024reporting,picard2021torch}.
Any study claiming that RNG quality interacts with optimiser choice (e.g., Adam benefits more from high-entropy shuffles than SGD) must therefore use multi-seed, multi-run designs.

\subsection{Explicit noise injection optimisers}

Beyond mini-batch noise, several methods inject randomness deliberately:
\begin{itemize}[leftmargin=*]
  \item \textbf{Stochastic Langevin / SGLD} add Gaussian noise scaled to mimic posterior sampling.
  \item \textbf{Gradient noise / parameter noise} exploration in reinforcement learning.
  \item \textbf{Entropy-regularised / noisy networks} policies.
\end{itemize}
For these algorithms, the distributional form of the injected noise is part of the method.
Substituting a non-Gaussian or correlated bitstream (via inverse-CDF transforms of biased bits) would constitute a meaningful ablation; we found little systematic work that varies \emph{bitstream quality} while keeping the intended noise law fixed through proper sampling transforms.
This is a concrete experimental gap (\S\ref{sec:gaps}).

\subsection{Metaheuristic and evolutionary optimisation}

Population-based methods consume large volumes of random numbers for initialisation and variation operators.

\subsubsection{Particle swarm optimisation (PSO)}

Bastos-Filho et al.~\cite{bastos2009pso} compare multiple RNGs (LCG variants and Marsaglia-type generators) across PSO topologies.
Their influential conclusion: PSO requires a \emph{minimum} RNG quality, but medium- and high-quality generators show no significant further gains.
GPU PSO studies similarly find that lightweight generators such as xorshift can be ``good enough'' when they clear basic statistical bars~\cite{bastos2010gpupso}.
Ma et al.~\cite{ma2014comparison} compare many library RNGs on GPU PSO with league-scoring protocols; differences exist but are function-dependent and often small among competent generators.

\subsubsection{Genetic algorithms and quasi-random initialisation}

Studies substituting pseudo- and \emph{quasi}-random sequences (Sobol, Halton, Niederreiter) for GA initial populations show that space-filling properties can matter more than fine differences among PRNGs, because early exploration depends on covering the search domain~\cite{cantu2008rngga,maaranen2004quasi}.
Reported rankings of generators are problem-dependent for constrained and unconstrained benchmarks---already a warning against universal claims~\cite{cantu2008rngga}.

\subsubsection{Differential evolution and hybrids}

Studies of DE/PSO hybrids under different PRNG quality levels report that weaker LCGs are only slightly worse than stronger generators in some office-solver settings, occasionally even improving convergence---consistent with the broader pattern that once catastrophic generators are excluded, effects are small and non-monotone~\cite{angelova2021de}.

\begin{remark}[Critical reading]
Metaheuristic RNG papers often vary generator algorithms without publishing NIST/TestU01 profiles of the exact streams consumed, and rarely power their comparisons for small effects.
Still, their repeated ``minimum quality threshold'' finding is strikingly consistent with Antunes et al.'s ML conclusion~\cite{antunes2025influence} and with Heese et al.'s null QRNG result~\cite{heese2024biased}.
\end{remark}

\subsection{Cross-cutting picture for optimisation}

\begin{table}[t]
\centering
\caption{How randomness interacts with major optimiser families.}
\label{tab:opt-map}
\begin{tabular}{@{}p{3.2cm}p{4.2cm}p{4.4cm}@{}}
\toprule
\textbf{Family} & \textbf{Primary randomness} & \textbf{Documented quality sensitivity} \\
\midrule
Mini-batch SGD & Data sampling / shuffle & Strong theory for noise \emph{structure}; little direct RNG-quality ablation \\
Adam / adaptive & Shuffle + moment estimation & Early-step variance issues~\cite{liu2020radam}; seed-sensitive rankings \\
Langevin-type & Explicit isotropic/anisotropic noise & Sensitive to noise law; bitstream quality rarely tested \\
PSO / DE / GA & Init + movement operators & Clear floor effects for weak RNGs; flat above medium quality~\cite{bastos2009pso,ma2014comparison} \\
Bayesian opt.\ / bandits & Acquisition sampling & Under-explored w.r.t.\ RNG quality \\
\bottomrule
\end{tabular}
\end{table}

Overall, optimisation research offers a sharper mechanistic story than ML substitution studies: \emph{structured} noise changes which solutions are stable.
But it almost never measures whether the RNG that implements sampling preserves that structure.
Bridging these literatures is a priority gap.

\section{Critical Synthesis}
\label{sec:critique}

We organise the critique around recurring failure modes rather than paper-by-paper commentary.

\subsection{Construct confusion: seeds, sources, and quality}

Many papers advertise ``effects of randomness'' while only sweeping seeds of a single PRNG.
Others swap QRNG for PRNG without publishing statistical profiles of either stream.
A third group publishes Crush scores for framework generators without measuring downstream task metrics.
These are three different scientific objects:
\begin{enumerate}[leftmargin=*]
  \item \textbf{Seed variance} --- sampling variability within a family;
  \item \textbf{Source identity} --- hardware versus algorithm brand;
  \item \textbf{Measured quality} --- test-battery / entropy / distinguisher scores.
\end{enumerate}
Only (iii) supports claims about ``quality of randomness.''
Source identity (ii) is a proxy that fails when a QRNG is biased~\cite{heese2024biased} or a PRNG is cryptographically strong.

\subsection{Contradictory ML substitution results}

Bird et al.~\cite{bird2019softcomputing} suggest meaningful QRNG benefits on some tasks; Heese et al.~\cite{heese2024biased} find null results even under hardware bias; Koivu et al.~\cite{koivu2022dropout} find signed effects that flip with dataset; Antunes et al.~\cite{antunes2025influence} argue for a quality floor with flatness above it.
These results are not necessarily logically inconsistent---effect modifiers (task difficulty, site of injection, stream length, preprocessing) can reconcile them---but the literature does not yet identify those modifiers with any reliability.

\subsection{Underpowered comparisons and selective reporting}

Effect sizes in RNG-substitution studies are often $<1$--$3\%$ absolute accuracy.
Reporting guidance for neural comparisons warns that small effects require large seed counts and non-parametric tests; otherwise conclusions are fragile or even adversarially selectable~\cite{akesson2024reporting}.
Few RNG-quality papers meet that bar.
Learning-curve plots without uncertainty bands, and ``QRNG wins'' headlines without multiplicity correction, remain common.

\subsection{Site confounding}

Initialisation, shuffle, dropout, and augmentation are frequently randomised together under one global seed.
Attributing a metric change to ``QRNG initialisation'' is then invalid if data order also changed.
Tooling nondeterminism (GPU reductions, cuDNN autotune) adds a further confound~\cite{zhuang2022randomness}.
Proper designs freeze all sites but one and disable nondeterministic kernels for confirmatory runs.

\subsection{Theory--practice mismatch on optimisation noise}

SGD theory cares about the \emph{covariance geometry} of gradient noise~\cite{zhu2019anisotropic,wu2022alignment}.
RNG test batteries care about uniformity and serial independence of bits.
No established reduction shows that passing BigCrush implies preserving SGD's alignment property under finite-precision shuffling of a real dataset.
Conversely, a weak LCG might still induce adequate mixing on a small dataset for few epochs.
The community lacks a \emph{task-relevant} quality metric for optimiser noise.

\subsection{Threshold phenomena versus romanticism about true randomness}

Across PSO~\cite{bastos2009pso}, framework PRNG studies~\cite{antunes2025influence}, and careful QRNG null results~\cite{heese2024biased}, a pragmatic pattern emerges: below a quality floor, algorithms break or degrade; above it, returns diminish quickly.
This undercuts both (a)~neglect of RNG choice (weak generators do hurt) and (b)~strong marketing claims that QRNG/TRNG must improve deep learning whenever substituted.
The interesting scientific question is the shape of the quality--utility curve---and whether ML sites differ in their floors.

\subsection{Publication and tooling biases}

Positive QRNG-advantage stories are more surprising than nulls; nulls may be under-published.
Industry stacks default to fast counter-based PRNGs without exposing bitstream diagnostics.
Energy and performance studies~\cite{antunes2024reproducibility} further push practitioners toward speed, not quality telemetry.
As a result, large random archives held by labs are rarely evaluated in ML loops with the same rigour used in cryptographic certification.

\section{Research Gaps}
\label{sec:gaps}

We state gaps as falsifiable absences in the literature.

\subsection{G1 --- No shared quality$\to$utility benchmark}

There is no community benchmark that (a)~publishes fixed random bit archives at multiple quality tiers, (b)~defines isolated stochastic-site hooks in reference trainers, and (c)~reports multi-seed metrics with pre-registered tests.
Existing papers each bring private streams and private tasks, preventing meta-analysis.

\subsection{G2 --- Missing isolation of stochastic sites}

Dropout-quality~\cite{koivu2022dropout} and initialisation-quality~\cite{bird2019softcomputing,heese2024biased} are studied separately; shuffle-quality and augmentation-quality almost not at all.
A full factorial (or at least one-site-at-a-time) design across S1--S4 (\S\ref{sec:background}) is absent.

\subsection{G3 --- Optimiser$\times$RNG interaction untested}

Despite rich theory on SGD versus Adam noise~\cite{zhou2020adamsgd,liu2020radam,wu2022alignment}, we found no controlled study that crosses optimiser family with measured bitstream quality of the shuffle/mask generator while holding all else fixed.
This is perhaps the highest-value gap for connecting the two literatures.

\subsection{G4 --- No task-relevant quality metric}

NIST/TestU01 failures do not map monotonically to accuracy deltas~\cite{antunes2025influence}.
Neural distinguishers~\cite{valle2024assessing} measure predictability, not optimiser utility.
Needed: metrics that predict downstream ML/optimiser sensitivity (e.g., estimated bias in mini-batch index distributions; autocorrelation at batch-scale lags; effective entropy consumed per epoch).

\subsection{G5 --- Scale mismatch}

Pilot studies use kilobytes to megabytes of random data.
Modern training can consume gigabytes of random bits across epochs (dropout masks alone).
Whether quality defects that appear only at Crush/BigCrush scales affect long training runs is unknown.
Large binary archives are uniquely suited to close this gap---yet remain under-used in published ML experiments.

\subsection{G6 --- Replication crisis on QRNG claims}

The Bird versus Heese disagreement~\cite{bird2019softcomputing,heese2024biased} is unresolved by direct, pre-registered replication on shared tasks with shared quality diagnostics.
Until then, QRNG-advantage claims should be treated as hypothesis-generating, not established.

\subsection{G7 --- Metaheuristics and deep learning are siloed}

PSO/GA/DE studies converged early on ``minimum quality suffices''~\cite{bastos2009pso,ma2014comparison}, while deep learning discussions rediscover similar ideas a decade later~\cite{antunes2025influence}.
Cross-domain synthesis---and tests of quasi-random versus pseudo-random \emph{mini-batch schedules}---are rare.

\subsection{G8 --- Reporting standards}

Papers seldom report: bytes consumed, wraparound policy, bit ordering, exact seeding maps, GPU determinism flags, or effect sizes with confidence intervals.
Without these, RNG-quality science cannot accumulate.

\subsection{Gap priority matrix}

\begin{table}[t]
\centering
\caption{Gap priorities for a bitstream-centred experimental programme.}
\label{tab:gap-priority}
\begin{tabular}{@{}clp{7.5cm}@{}}
\toprule
\textbf{ID} & \textbf{Priority} & \textbf{Rationale} \\
\midrule
G1 & High & Enables cumulative science \\
G2 & High & Removes attribution confounds \\
G3 & High & Links RNG quality to optimiser theory \\
G4 & Medium & Needs iterative metric design \\
G5 & High & Motivates large archives \\
G6 & Medium & Clarifies QRNG narrative \\
G7 & Medium & Broadens optimisation coverage \\
G8 & High & Cheap to adopt; high leverage \\
\bottomrule
\end{tabular}
\end{table}

\section{Proposed Methodology}
\label{sec:methodology}

This section converts the gaps of \S\ref{sec:gaps} into an operational methodology.
The design is intentionally narrow: it prioritises \emph{causal attribution} of bitstream quality over breadth of tasks.

\subsection{Design principles}

\begin{enumerate}[leftmargin=*]
  \item \textbf{Measure quality, do not assume it.}
  Each bitstream receives Stage~A scores (cheap statistical suite; NIST/TestU01 when data volume allows; optional neural distinguisher~\cite{valle2024assessing}) and a tier label $T0$--$T4$.

  \item \textbf{Isolate one stochastic site.}
  Following the control logic of Summers and Dinneen~\cite{summers2021nondeterminism}, when site $i$ is under test, all sites $j\neq i$ are frozen (fixed PRNG seeds), the train/test split is fixed, and confirmatory runs use deterministic compute kernels where possible~\cite{zhuang2022randomness}.

  \item \textbf{Power the comparison.}
  Use a predefined run index set, non-parametric tests (Brunner--Munzel), effect sizes (Cliff's $\delta$ / stochastic superiority), and $n\geq 25$ when expected gaps are below $2.5\%$~\cite{akesson2024reporting}.

  \item \textbf{Forbid silent wraparound.}
  Confirmatory runs must not recycle a short buffer without logging; byte consumption is an audit field.
\end{enumerate}

\subsection{Independent and dependent variables}

\begin{table}[t]
\centering
\caption{Core methodological variables.}
\label{tab:vars}
\begin{tabular}{@{}p{3.4cm}p{8.4cm}@{}}
\toprule
\textbf{Variable} & \textbf{Operationalisation} \\
\midrule
Bitstream quality & Stage~A scores + tier $T0$--$T4$; source identity recorded but not used as the sole factor \\
Stochastic site & Dropout masks; data shuffle; weight init; (later) augmentation \\
Optimiser & SGD vs.\ Adam with fixed hyperparameters within a comparison cell \\
Task & Fashion-{MNIST} MLP (pilot); CIFAR-10 / noisy tabular (main) \\
Outcomes & Test accuracy / F1, train--test gap, curve AUC, run-to-run IQR \\
\bottomrule
\end{tabular}
\end{table}

\subsection{Bitstream consumption interface}

We require a \texttt{BitstreamRNG} that memory-maps or streams a binary archive and exposes uniform draws, Bernoulli masks, shuffles, and Gaussian transforms derived from those uniforms.
Quality is therefore a property of the \emph{file}, not of an opaque framework seed.
Starting offsets / shards provide replication structure without reseeding a classical PRNG that would bypass the archive.

\subsection{Priority experiments}

\begin{enumerate}[leftmargin=*]
  \item \textbf{Negative-control gate (dropout).}
  Biased bits ($p=0.6$) versus OS CSPRNG drive dropout only; $n=10$.
  Pipeline validity requires a detectable difference in accuracy, gap, or dispersion.

  \item \textbf{Quality ladder (dropout).}
  Weak LCG, MT19937-class stream, CSPRNG, and a large-archive shard; same isolation; escalate to $n\geq 25$ if gaps are small.

  \item \textbf{Optimiser interaction (shuffle).}
  Cross shuffle-bitstream quality with $\{\mathrm{SGD},\mathrm{Adam}\}$ to test whether adaptive methods amplify or damp bitstream defects in gradient noise---addressing gap~G3.

  \item \textbf{Initialisation replication cell.}
  Init-only ablation in the style of Bird et al.~\cite{bird2019softcomputing} and Heese et al.~\cite{heese2024biased}, but with published Stage~A scores attached to every source.
\end{enumerate}

\subsection{What this methodology deliberately defers}

Full Crush batteries on undersized pilots; joint multi-site archive injection; large hyperparameter sweeps; and metaheuristic suites are deferred until the dropout negative-control gate passes.
Quasi-random mini-batch schedules and neural distinguishers are optional enrichments, not blockers.

\subsection{Claim policy}

A result may be stated as a \emph{quality effect} only if Stage~A scores differ across compared sources, a single site was ablated, wraparound policy is stated, and the statistical protocol above is followed.
Null results among $T2$--$T4$ sources are considered scientifically informative when adequately powered.

\section{Research Agenda}
\label{sec:agenda}

We outline a programme tailored to laboratories that hold large binary random archives and wish to test quality effects in AI/ML and optimisation.
Details of file formats and pilot generators appear in the accompanying project documentation; here we state the scientific protocol.

\subsection{Two-stage design}

\paragraph{Stage A --- Characterise.}
For each archive (and for negative/positive controls: biased bits, weak LCG, MT19937, OS CSPRNG):
\begin{itemize}[leftmargin=*]
  \item run a cheap statistical suite (frequency, runs, serial correlation, byte $\chi^2$, Shannon entropy);
  \item run NIST SP~800-22 and/or TestU01 SmallCrush/Crush as data volume allows;
  \item optionally train a neural distinguisher against a golden reference~\cite{valle2024assessing};
  \item assign a tier label $T0$--$T4$ from fail-basic to golden-archive.
\end{itemize}

\paragraph{Stage B --- Inject.}
Implement a \texttt{BitstreamRNG} that consumes archive bytes without wraparound in confirmatory runs.
Ablate one stochastic site at a time (recommended order: dropout $\rightarrow$ shuffle $\rightarrow$ initialisation $\rightarrow$ augmentation).
Cross with optimiser family $\{$\textrm{SGD}, \textrm{SGD+momentum}, \textrm{Adam}$\}$ on at least one vision and one tabular/noisy task.

\subsection{Minimal decisive experiment}

\begin{enumerate}[leftmargin=*]
  \item \textbf{Negative control check.} Biased ($p=0.6$) bits versus OS CSPRNG on dropout masks for an MLP on Fashion-MNIST; $n=10$ runs.
  If no difference in accuracy or train--test gap, debug injection before claiming nulls among strong sources.
  \item \textbf{Quality ladder.} Weak LCG, MT19937, OS CSPRNG, and one large archive shard; same protocol; $n\geq 25$ if observed gaps $<2.5\%$~\cite{akesson2024reporting}.
  \item \textbf{Optimiser interaction.} Repeat the ladder under SGD and Adam with fixed hyperparameters other than the bitstream-driven site.
\end{enumerate}

\subsection{Statistical protocol}

Predefine seeds $\{0,\ldots,n-1\}$ for any residual PRNG needs; use Brunner--Munzel or Mann--Whitney tests; report Cliff's $\delta$ or Cohen's $d$; correct for multiplicity across sources within a site (Holm).
Publish full result tables and bitstream manifests (hashes, bytes consumed).

\subsection{Success criteria for the research programme}

The programme succeeds scientifically even if most strong sources are indistinguishable:
\begin{itemize}[leftmargin=*]
  \item a measured quality--utility curve (possibly flat above a floor);
  \item site and optimiser interaction maps;
  \item open artefacts enabling replication (G1, G8).
\end{itemize}
A positive finding (tier predicts accuracy/stability) would justify operational use of high-entropy archives in training; a well-powered null among $T2$--$T4$ sources would equally inform practitioners to invest elsewhere (data, hyperparameters) once a quality floor is met.

\subsection{Conclusion}

Randomness is indispensable to ML and optimisation, yet ``quality of randomness'' remains an underspecified lever.
The literature shows clear harm from weak generators, suggestive but conflicted benefits from QRNG/TRNG substitutions, and a mature theory of \emph{structured} optimiser noise that has not been connected to bitstream diagnostics.
Closing that gap requires experiments that measure quality, isolate stochastic sites, cross optimiser families, and power their statistics---precisely the setting for which large binary random archives are an asset rather than a curiosity.


\section*{Acknowledgements}
This review was prepared as part of the Randomness project experimental programme.
% Add funding / collaborators as appropriate.

\bibliography{refs}

\end{document}
